\chapter{Cholesteatoma Surgery}

\section*{Introduction \& Clinical Indications}
Cholesteatoma surgery is a specialized surgical procedure aimed at the complete removal of a cholesteatoma, which is an abnormal growth of skin cells in the middle ear and mastoid process. This destructive lesion, often described as a "skin cyst" in an inappropriate location, can lead to significant complications if left untreated. Cholesteatomas can erode surrounding bone structures, including the ossicles, facial nerve canal, and even the inner ear, resulting in hearing loss and potential intracranial complications such as meningitis or brain abscess. The surgery is critical not only for eradicating the cholesteatoma but also for preserving hearing and preventing further morbidity associated with the disease. The procedure is typically performed under general anesthesia and involves meticulous dissection to ensure complete excision of the cholesteatoma matrix while safeguarding adjacent vital structures.

\begin{itemize}
  \item Mastoidectomy (often combined with tympanoplasty or ossiculoplasty).
  \item Tympanomastoidectomy.
  \item Cholesteatoma excision.
  \item Canal wall up (intact canal wall) mastoidectomy.
  \item Canal wall down (open cavity) mastoidectomy.
  \item Modified radical mastoidectomy.
  \item Radical mastoidectomy (for extensive disease).
\end{itemize}

The primary surgical principle in cholesteatoma surgery is the complete excision of the cholesteatoma sac and its keratin debris from the middle ear and mastoid. The cholesteatoma is characterized by continuous desquamation of keratin, leading to its expansion and subsequent pressure necrosis of surrounding tissues. The rationale for surgical intervention is to halt this destructive process, prevent further hearing loss, and eliminate the risk of serious complications such as intracranial infections. The technical goals include achieving a safe, dry, and disease-free ear, preserving or reconstructing the sound-conducting mechanism (ossicular chain), and ensuring a stable tympanic membrane. The choice between an intact canal wall (canal wall up) or an open cavity (canal wall down) approach depends on the extent of disease, the surgeon's preference, and the patient's compliance with follow-up care.

The management of cholesteatoma has evolved significantly over the past century. Early attempts in the 19th century focused on drainage and debridement, often resulting in high morbidity. Pioneers such as Josef Gruber and Hermann Schwartze, in 1873, introduced mastoidectomy techniques that involved opening the mastoid cortex to drain infected material. The introduction of the operating microscope by Carl Olof Nyl\'{e}n in 1922 revolutionized otologic surgery, allowing for more precise dissection and preservation of vital structures. Surgeons like John Shea and Michael Portmann advanced techniques aimed at preserving the ear canal, leading to the development of the intact canal wall mastoidectomy. However, challenges with residual disease prompted refinements in canal wall down procedures, such as the modified radical mastoidectomy, which provides better visualization and facilitates easier long-term care. Today, cholesteatoma surgery often employs a staged approach, with initial removal of the disease followed by second-look surgeries to assess for residual cholesteatoma and perform ossicular reconstruction.

Cholesteatoma surgery is indicated for patients presenting with confirmed or highly suspected cholesteatoma, typically manifesting with the following:

\begin{itemize}
  \item Chronic, persistent, and often foul-smelling otorrhea (ear discharge) unresponsive to medical therapy.
  \item Progressive conductive or mixed hearing loss due to ossicular chain erosion or tympanic membrane perforation.
  \item Otoscopic visualization of a retraction pocket, keratin debris, or a pearly white mass behind an intact or perforated tympanic membrane.
  \item Vertigo, dizziness, or imbalance, suggesting erosion into the bony labyrinth (e.g., lateral semicircular canal fistula).
  \item Facial nerve weakness or paralysis, indicating erosion of the fallopian canal and compression or irritation of the nerve.
  \item Headache, fever, or other signs of intracranial complications such as meningitis, brain abscess, or sigmoid sinus thrombosis.
  \item Radiographic evidence on high-resolution computed tomography (CT) of the temporal bone showing a soft tissue mass with associated bone erosion (e.g., scutum erosion, ossicular erosion, tegmen erosion).
  \item Recurrent acute otitis media in the presence of a deep retraction pocket or suspected cholesteatoma.
\end{itemize}

Cholesteatoma surgery is considered the primary and definitive treatment for established cholesteatoma. Unlike other ear conditions that may respond to medical management, no pharmacological therapy can eradicate the epithelial sac of a cholesteatoma. Observation alone carries significant risks of progressive bone destruction, irreversible hearing loss, facial nerve injury, and potentially fatal intracranial complications. Therefore, once a cholesteatoma is diagnosed, surgical extirpation is the first-line and often urgent intervention. In select cases of small, quiescent, and non-progressive retraction pockets, a period of watchful waiting with regular cleaning may be considered, but this is an exception requiring strict patient compliance. Surgery may also be staged, with an initial procedure to remove the bulk of the disease and a second-look operation to check for residual disease and perform definitive ossicular reconstruction.

\section*{Relevant Surgical Anatomy}
The surgical management of cholesteatoma primarily involves the middle ear and mastoid air cell system, both located within the temporal bone. Key anatomical structures include:

- **Tympanic Membrane**: The eardrum, which forms the lateral boundary of the middle ear.
- **Ossicular Chain**: Comprising the malleus, incus, and stapes, these small bones transmit sound vibrations to the inner ear.
- **Facial Nerve (Cranial Nerve VII)**: This nerve traverses the temporal bone within the fallopian canal, making it vulnerable during surgical dissection.
- **Chorda Tympani Nerve**: A branch of the facial nerve that provides taste sensation to the anterior two-thirds of the tongue, crossing the middle ear cavity.
- **Inner Ear Structures**: Including the cochlea and semicircular canals, which are encased in bone and form the bony labyrinth.
- **Tegmen Tympani and Tegmen Mastoideum**: These structures form the superior bony roof separating the middle ear and mastoid from the middle cranial fossa.
- **Sigmoid Sinus**: A major venous channel located posteriorly within the mastoid.

Cholesteatomas can extend into the mastoid air cells, eroding bone and placing adjacent vital structures at risk. Understanding these anatomical relationships is crucial for successful surgical intervention.

\section*{Preoperative Workup \& Preparation}
A comprehensive preoperative workup is essential to assess the extent of disease, identify potential complications, and optimize patient safety.

\begin{itemize}
  \item \textbf{Clinical Evaluation:} Detailed otologic history and physical examination, including otomicroscopy, to assess the tympanic membrane, ear canal, and presence of discharge.
  
  \item \textbf{Audiometry:} Pure tone audiometry and speech audiometry are performed to quantify the type and degree of hearing loss.
  
  \item \textbf{Imaging:} A high-resolution computed tomography (CT) scan of the temporal bone is mandatory. It delineates the extent of bone erosion, ossicular status, integrity of the tegmen, presence of labyrinthine fistula, and facial nerve canal involvement. Magnetic resonance imaging (MRI) may be used in select cases to differentiate cholesteatoma from granulation tissue or to evaluate for intracranial extension.
  
  \item \textbf{Infection Control:} Any active ear infection or purulent discharge is treated with appropriate topical or oral antibiotics and meticulous ear cleaning to create a dry, sterile surgical field, which improves graft success rates.
  
  \item \textbf{Medical Clearance:} Patients undergo routine preoperative medical evaluation, including blood tests (CBC, electrolytes, coagulation profile), urinalysis, and an electrocardiogram (ECG). Patients with significant comorbidities may require cardiac or pulmonary clearance.
  
  \item \textbf{Medication Adjustment:} Anticoagulant and antiplatelet medications are typically discontinued or adjusted according to institutional protocols to minimize bleeding risk.
  
  \item \textbf{NPO Status:} Patients are instructed to be NPO (nil per os) for at least 6-8 hours prior to general anesthesia.
  
  \item \textbf{Informed Consent:} Detailed discussion with the patient regarding the nature of the disease, the surgical procedure, potential risks, benefits, alternative treatments, and the need for long-term follow-up.
  
  \item \textbf{Prophylactic Antibiotics:} Intravenous prophylactic antibiotics are administered prior to incision.
\end{itemize}

While cholesteatoma surgery is often imperative, certain conditions may contraindicate or necessitate precautions for the procedure.

\textbf{Factors requiring optimization or a modified approach:}
\begin{itemize}
  \item Severe, uncorrectable systemic illness may require optimization, a different anesthetic plan, or urgent risk-benefit review.
  
  \item Uncontrolled, severe coagulopathy that cannot be corrected.
\end{itemize}

\textbf{Important precautions:}
\begin{itemize}
  \item A small, quiescent, non-progressive retraction pocket that can be reliably observed and cleaned medically, especially in an elderly or frail patient.
  
  \item Active drainage or infection should be treated and the ear cleaned when possible, but infection does not eliminate the need for surgery when destructive disease is present.
  
  \item A contralateral only-hearing ear requires careful counseling, hearing-preservation planning, and discussion of the risk of further hearing loss.
  
  \item Patient non-compliance with follow-up, which is crucial for monitoring recurrence, especially with intact canal wall techniques.
  
  \item Significant medical comorbidities that increase the risk of anesthesia or surgery beyond the immediate threat posed by the cholesteatoma.
  
  \item Imaging evidence of erosion into the lateral semicircular canal (labyrinthine fistula) or the fallopian canal (facial nerve canal), requiring meticulous dissection and potentially a staged approach to minimize inner ear trauma or facial nerve injury.
  
  \item Extensive intracranial extension, which may necessitate a combined neurosurgical-otologic approach.
\end{itemize}

\section*{Surgical Technique \& Procedure}
Cholesteatoma surgery is performed under general anesthesia, typically with the patient in a supine position and the head turned to the contralateral side to expose the affected ear. Facial nerve monitoring is routinely employed to protect the nerve during dissection.

\begin{itemize}
  \item \textbf{Anesthesia \& Positioning:} General endotracheal anesthesia is administered. The patient is positioned supine with the head turned, and the surgical field is prepped and draped. Facial nerve monitoring electrodes are placed.
  
  \item \textbf{Incision:} The approach varies based on disease extent. A postauricular (behind the ear) incision is common for mastoidectomy, allowing wide exposure. A transcanal or endaural incision may be used for limited disease.
  
  \item \textbf{Exposure:} The skin and soft tissues are elevated to expose the mastoid cortex and external auditory canal. For postauricular approaches, the temporalis fascia is often harvested for use as a graft.
  
  \item \textbf{Mastoidectomy:} Using a high-speed surgical drill with continuous irrigation, the mastoid air cells are systematically opened. The extent of mastoidectomy (e.g., cortical, extended, radical) depends on the disease spread. The cholesteatoma sac is identified and meticulously dissected from surrounding vital structures.
  
  \item \textbf{Cholesteatoma Removal:} Under high magnification, the cholesteatoma matrix is carefully separated from the facial nerve, ossicles, labyrinth, and dura. Complete removal is paramount to prevent recurrence. If a labyrinthine fistula is present, the matrix over the fistula may be left in situ and covered with bone dust or fascia to avoid inner ear trauma.
  
  \item \textbf{Ossicular Reconstruction (Ossiculoplasty):} After disease removal, the ossicular chain is inspected. If eroded, it is reconstructed using autologous bone (e.g., incus, malleus head), cartilage, or synthetic prostheses (e.g., titanium prostheses) to restore sound conduction.
  
  \item \textbf{Tympanic Membrane Repair (Tympanoplasty):} A perforation in the tympanic membrane is repaired using a graft, most commonly temporalis fascia or tragal cartilage, placed either underlay or overlay.
  
  \item \textbf{Closure:} The middle ear and mastoid cavity may be packed with absorbable gelatin sponge or other materials. The ear canal is packed with antibiotic-impregnated gauze. The incision is closed in layers with sutures, and a mastoid dressing is applied. In canal wall down procedures, the mastoid cavity is exteriorized into the ear canal, and no external dressing is typically used over the cavity itself, allowing for postoperative cleaning.
\end{itemize}

\section*{Complications \& Risks}
Cholesteatoma surgery, while highly effective, carries inherent risks, both general to any surgery and specific to otologic procedures.

\textbf{General Surgical Complications:}
\begin{itemize}
  \item Bleeding (hematoma formation, intraoperative hemorrhage).
  
  \item Infection (wound infection, mastoiditis).
  
  \item Adverse reaction to anesthesia (nausea, vomiting, allergic reactions, cardiac events).
  
  \item Pain.
  
  \item Scarring.
\end{itemize}

\textbf{Procedure-Specific Complications:}
\begin{itemize}
  \item \textbf{Hearing Loss:} Conductive hearing loss may persist or worsen if ossicular reconstruction fails. Sensorineural hearing loss can occur due to inner ear trauma (e.g., drilling near the labyrinth, labyrinthine fistula manipulation).
  
  \item \textbf{Facial Nerve Injury:} Temporary or permanent facial weakness or paralysis due to direct trauma, thermal injury from drilling, or compression. This is a rare but devastating complication.
  
  \item \textbf{Dizziness and Vertigo:} Postoperative vertigo is common, especially if a labyrinthine fistula was present or manipulated. Persistent vertigo can occur due to inner ear damage.
  
  \item \textbf{Tinnitus:} New or worsened ringing in the ear.
  
  \item \textbf{Recurrence or Residual Cholesteatoma:} Despite meticulous surgery, microscopic remnants of epithelium can lead to recurrence, necessitating further surgery.
  
  \item \textbf{Taste Disturbance:} Injury to the chorda tympani nerve, which traverses the middle ear, can cause altered taste sensation on the anterior two-thirds of the tongue.
  
  \item \textbf{Tympanic Membrane Perforation/Graft Failure:} The eardrum graft may fail to take, resulting in a persistent perforation.
  
  \item \textbf{Mastoid Cavity Problems (Canal Wall Down):} An open mastoid cavity requires lifelong cleaning and can be prone to discharge, infection, or difficulty with water exposure.
  
  \item \textbf{Intracranial Complications:} Extremely rare but serious risks include meningitis, brain abscess, cerebrospinal fluid (CSF) leak, or sigmoid sinus thrombosis if the dura or major venous sinuses are breached.
\end{itemize}

\section*{Postoperative Care \& Recovery}
The postoperative recovery from cholesteatoma surgery varies depending on the extent of the procedure and the individual patient's healing capacity. 

Immediately after surgery, patients are monitored in the Post-Anesthesia Care Unit (PACU) for recovery from general anesthesia, pain control, and vital sign stability. Most patients experience mild to moderate ear pain, which is managed with oral analgesics. Nausea and dizziness are common, especially in the first 24-48 hours, and antiemetics are prescribed as needed. A mastoid dressing is typically in place, and patients are instructed to keep the ear dry and avoid strenuous activities. 

For the first 1-2 weeks, patients are advised to avoid heavy lifting, straining, nose blowing, and any activities that increase intracranial pressure, to protect the healing graft and prevent bleeding. Diet is usually advanced as tolerated. The first postoperative visit typically occurs within 1-3 weeks, where the external dressing and ear canal packing are removed, and the ear is inspected under the microscope. Patients are usually able to return to light activities and work within 2-4 weeks, but full recovery, including resolution of swelling and complete graft healing, can take several months. Water precautions (e.g., earplugs for showering) are usually maintained until the surgeon confirms complete healing of the eardrum. Long-term follow-up, often for many years, is crucial due to the risk of cholesteatoma recurrence.

During surgery, the surgeon meticulously documents the extent of the cholesteatoma, its involvement with the ossicles, facial nerve, labyrinth, and tegmen, and the type of reconstruction performed. Key findings include the size and location of the cholesteatoma matrix, the presence and extent of ossicular erosion, any labyrinthine fistula, and the condition of the middle ear mucosa. The removed tissue, consisting of the cholesteatoma sac and keratin debris, is sent for histopathological examination. The pathology report confirms the diagnosis of cholesteatoma by identifying stratified squamous epithelium with keratinizing features within the specimen. It also notes the presence of inflammatory cells, granulation tissue, and any associated bone fragments. This report is crucial for confirming the nature of the disease and guiding further management, especially in cases where the intraoperative diagnosis was uncertain.

A normal and successful postoperative course following cholesteatoma surgery is characterized by a gradual resolution of pain and discomfort, followed by complete healing of the surgical site. The ear should become dry, with no persistent discharge. The reconstructed tympanic membrane graft should be intact and well-vascularized. Hearing, while not always fully restored to normal, should be preserved or improved, particularly the conductive component. The patient should experience no new or worsening facial weakness, vertigo, or tinnitus. For patients undergoing canal wall down procedures, the mastoid cavity should remain clean and dry, either self-cleaning or easily maintained with occasional clinic cleaning. Long-term, the absence of recurrent or residual cholesteatoma on clinical examination and imaging (if performed) signifies a successful outcome, allowing the patient to return to normal activities without the threat of progressive disease.

Postoperative care for cholesteatoma surgery is critical for monitoring healing, detecting complications, and ensuring long-term disease control.

\begin{itemize}
  \item \textbf{Initial Postoperative Visits:} Typically within 1-3 weeks for removal of external dressings and ear canal packing, followed by microscopic examination of the ear.
  
  \item \textbf{Regular Microscopic Surveillance:} Patients require periodic follow-up visits, often every 3-6 months initially, then annually for many years. This allows the surgeon to inspect the ear canal, tympanic membrane, and mastoid cavity (if open) for any signs of recurrence or residual disease.
  
  \item \textbf{Audiometry:} Repeat audiograms are performed several months postoperatively to assess the final hearing outcome after healing is complete.
  
  \item \textbf{Imaging:} In cases of intact canal wall mastoidectomy, a follow-up high-resolution CT scan or diffusion-weighted MRI may be performed 6-12 months postoperatively to screen for residual cholesteatoma, especially if a second-look surgery is not planned.
  
  \item \textbf{Water Precautions:} Strict avoidance of water in the ear until the tympanic membrane graft is fully healed, typically 6-12 weeks.
  
  \item \textbf{Activity Restrictions:} Avoidance of strenuous activities, heavy lifting, and activities that cause pressure changes (e.g., diving, flying without pressure equalization) for several weeks to months.
  
  \item \textbf{Warning Signs:} Patients are educated to contact their surgeon immediately if they experience new or worsening pain, fever, persistent or foul-smelling ear discharge, sudden hearing loss, facial weakness, or severe vertigo. These symptoms may indicate infection, graft failure, or a serious complication.
  
  \item \textbf{Revision Surgery:} If residual or recurrent cholesteatoma is detected, or if the eardrum graft fails, revision surgery may be necessary.
\end{itemize}

\section*{Outcomes \& Prognosis}
Cholesteatoma is a relatively uncommon condition, with an estimated incidence ranging from 3 to 15 cases per 100,000 children and approximately 9 per 100,000 adults annually. Acquired cholesteatoma, which develops secondary to chronic otitis media and eustachian tube dysfunction leading to tympanic membrane retraction, accounts for the vast majority of cases. Congenital cholesteatoma, arising from embryonic epithelial rests, is rarer and typically presents as a pearly white mass behind an intact tympanic membrane. There is a slight male predominance in incidence. While it can occur at any age, acquired cholesteatoma is frequently diagnosed in children and young adults, often presenting with chronic ear discharge and progressive hearing loss. Bilateral disease is uncommon but can occur, particularly in patients with underlying conditions affecting eustachian tube function. Genetic predisposition may play a role in a subset of patients, as a positive family history is sometimes observed. The morbidity of cholesteatoma primarily stems from progressive hearing loss and the potential for severe, albeit rare, intracranial complications such as meningitis, brain abscess, and lateral sinus thrombosis, which underscore the importance of timely surgical intervention.

The prognosis for cholesteatoma following complete surgical removal is generally favorable, with the primary goals being a safe, dry, and disease-free ear, along with preservation or improvement of hearing. Success rates for achieving a dry, safe ear are high, often exceeding 90\%. However, the risk of residual or recurrent cholesteatoma remains a significant concern, with reported rates varying widely from 5\% to 30\% depending on the extent of the initial disease, the surgical technique employed (intact canal wall vs. canal wall down), and the duration and rigor of follow-up. Intact canal wall techniques tend to have a higher rate of recurrence but offer a more anatomically normal ear, while canal wall down procedures, by exteriorizing the mastoid, generally have lower recurrence rates but result in a cavity requiring lifelong care. Hearing outcomes are variable and depend heavily on the preoperative status of the ossicular chain and middle ear mucosa, as well as the success of ossicular reconstruction. Prognostic factors for poorer outcomes include extensive disease with erosion into the labyrinth or facial nerve canal, significant preoperative hearing loss, and the presence of a labyrinthine fistula. Despite these challenges, skilled surgery can often control the disease, prevent life-threatening complications, and significantly improve the patient's quality of life. Long-term surveillance is paramount to detect and manage any recurrence promptly.

\section*{References}
\begin{enumerate}
  \item MedlinePlus. Cholesteatoma surgery. U.S. National Library of Medicine; 2025.
  
  \item Cummings Otolaryngology - Head and Neck Surgery. 8th ed. Elsevier; 2021.
  
  \item Sabiston Textbook of Surgery: The Biological Basis of Modern Surgical Practice. 21st ed. Elsevier; 2021.
  
  \item American Academy of Otolaryngology - Head and Neck Surgery. Clinical Practice Guideline: Otitis Media with Effusion (Update). Otolaryngology--Head and Neck Surgery. 2016;154(1 Suppl):S1-S41.
\end{enumerate}

\clearpage
